\documentclass{article}
\usepackage{charter}
\usepackage[margin=1in]{geometry}
\usepackage{fancyhdr,graphicx}
\usepackage{xspace}
\usepackage[inline]{enumitem}
\usepackage{hyperref}
\hypersetup{
	colorlinks=true,
	linkcolor=blue,
	filecolor=magenta,      
	urlcolor=blue,
	pdftitle={Overleaf Example},
	pdfpagemode=FullScreen,
}

\fancypagestyle{firstpage}{%
	\fancyhf{} % Clear header/footer
	\renewcommand{\headrulewidth}{0pt}%
	\renewcommand{\footrulewidth}{1pt}%
	% Add more detail here if needed
}
\fancypagestyle{otherpages}{%
	\fancyhf{}% Clear header/footer
	\renewcommand{\headrulewidth}{1pt}%
	\renewcommand{\footrulewidth}{1pt}%
	% Add more detail here if needed
}

\AtBeginDocument{\thispagestyle{firstpage}}
\pagestyle{otherpages}

\setlength{\parindent}{0pt}
\setlength{\parskip}{1ex}

\begin{document}
	
	\vspace*{\dimexpr-\headsep-\headheight-1pt}
	
	% \includegraphics[width=2.75in]{logo.png}
	
	\rule{\linewidth}{1pt}
	
	\bigskip
	
	
	%----------------------------------------------------------------------------------------
	%	YOU ONLY NEED TO FILL IN THE FOLLOWING LATEX MACROS
	%----------------------------------------------------------------------------------------
	
	% Title of the paper
	\newcommand{\PaperTitle}{Magnetic damping of standing waves in rectangular geometries}
	
	% Full name of the authors
	\newcommand{\AuthorOne}{Pranav Hegde}
	\newcommand{\AuthorTwo}{Thomas Gundrum}
	\newcommand{\AuthorThree}{Gerrit Maik Horstmann}
	
	% Institution
	\newcommand{\Institution}{Helmholtz-Zentrum Dresden-Rossendorf\xspace}
	\newcommand{\InstitutionAddress}{Bautzner Landstr. 400\xspace}
	\newcommand{\City}{Dresden, 01328\xspace}
	\newcommand{\Country}{Germany\xspace}
	
	% Full name of the Journal's Editor in Chief
	\newcommand{\EditorInChief}{Prof. Beverley J. McKeon\xspace}
	
	% Journal
	\newcommand{\Journal}{Physical Review Fluids\xspace}
	
	
	
	% % Body
	% \newcommand{\Contribution}{X\xspace}
	% \newcommand{\Problem}{Y\xspace}
	% \newcommand{\Solution}{Z\xspace}
	% \newcommand{\Evaluation}{D\xspace}
	% \newcommand{\NoveltyClaim}{E\xspace}
	
	%% List of contributions
	%\newcommand{\ContributionOne}{Contribution A.\xspace}
	%\newcommand{\ContributionTwo}{Contribution B.\xspace}
	%\newcommand{\ContributionThree}{Contribution C.\xspace}
	%\newcommand{\ContributionFour}{Contribution D.\xspace}
	
	% Suggested editors
	% \newcommand{\EditorOne}{Editor \#1\xspace}
	% \newcommand{\EditorOneExpertise}{A\xspace}
	
	% \newcommand{\EditorTwo}{Editor \#2\xspace}
	% \newcommand{\EditorTwoExpertise}{B\xspace}
	
	% \newcommand{\EditorThree}{Editor \#3\xspace}
	% \newcommand{\EditorThreeExpertise}{C\xspace}
	
	%----------------------------------------------------------------------------------------
	%	YOUR NAME AND CONTACT INFORMATION
	%----------------------------------------------------------------------------------------
	
	\hfill
	\begin{tabular}{ l @{} }
		\today \\[12pt] % Date
		\Institution\\
		\InstitutionAddress\\
		\City, \Country
	\end{tabular}
	
	
	%----------------------------------------------------------------------------------------
	%	LETTER CONTENT
	%----------------------------------------------------------------------------------------
	
	\bigskip
	
	Dear \EditorInChief,\\Chief Editor of \textit{\Journal}
	
	\bigskip
	
	We wish to submit an original journal article titled “\PaperTitle” in the \Journal. This study primarily aims to calculate the magnetic damping rate of a free-surface standing wave of a conducting fluid within a rectangular domain.     
	
	In summary, our contributions are the following:
	
	\begin{itemize}
		\item We formulate the magnetic damping rate of a surface gravity wave in a rectangular geometry filled with an electrically conducting fluid, which is subjected to a vertical magnetic field. 
		
		\item Previous studies only considered \begin{enumerate*}[label={\roman*)}] 
			\item only one single type of wave mode $(1,0)$, 
			\item shallow water approximation,
			\item and uniform electrical boundary conditions (either all insulating or all conducting).
		\end{enumerate*}
		The present study considered arbitrary wave modes with mixed electrical boundary conditions, disregarding the shallow water approximation.  
		
		\item Furthermore, we have also formulated a new expression for the Hartmann and Shercliff boundary layer in a single-layer system, and the Hartmann boundary layer formed at the interface in a two-layer system with large electrical conductivity difference (quasi one-layer system). 
	\end{itemize}
	
	We have also validated a part of the theory with a simple lab-scale sloshing experiment. Lastly, we provide a supplementary material that details all the analytical derivation along with a Python module to aid the readers in reproducing and applying our theoretical results. 
	
	I confirm that this paper, has not been published elsewhere. All the authors have approved the manuscript and agree with its submission to the \Journal.
	
	\bigskip
	Kind regards,
	
	\vspace{10pt}
	\AuthorOne,
	
	On behalf of \AuthorTwo \hspace{0.3mm} and\,  \AuthorThree
	
\end{document}